%% kapa: exam=1 round=32 year=2021 period=2 form=A subject=accounting subject_ko=회계학 seq=16 number=56
%% source: Q-Net 감정평가사 기출문제 (cst00309) artlSeq=5208397 file=extracted/1008/5208397/A형/2021년도 제32회 감정평가사 1차 2교시 A형.pdf
%% conversion: kordoc 4.24.51 -> tools/gen1.py (auto); review=proofread-sonnet; flags=none
\begin{question}{56}
(주)감평은 20×1년 초 기계장치(내용연수 3년, 잔존가치 ￦0, 정액법 상각)를 구입과 동시에 무이자부 약속어음(액면금액 ￦300,000, 3년 만기, 매년 말 ￦100,000 균등상환)을 발행하여 지급하였다. 이 거래 당시 (주)감평이 발행한 어음의 유효이자율은 연 12\%이다. 기계장치에 대해 원가모형을 적용하고, 당해 차입원가는 자본화대상에 해당하지 않는다. 20×1년 (주)감평이 인식할 비용은? (단, 12\%, 3기간의 연금현가계수는 2.40183이고, 계산금액은 소수점 첫째자리에서 반올림하며, 단수차이로 인한 오차가 있으면 가장 근사치를 선택한다.)
\begin{choices}
  \item ￦59,817
  \item ￦80,061
  \item ￦88,639
  \item ￦108,883
  \item ￦128,822
\end{choices}
\end{question}
\answer{4}
